\subsection{离散测度}

命题 14. --- 为使测度 \(\mu\) （局部紧空间 \(X\) 上的）是离散测度（§1, No.~3, 例 I），必须且只须 \(\operatorname{Supp}(\mu)\) 是 \(X\) 的一个离散闭子空间。

设 \(\mu\) 是 X 上的一个离散测度，由质量 \(h(x) \neq 0\) （放在 X 的离散闭子空间 N 的各点 \(x\) 处）所定义；我们证明 \(\operatorname{Supp}(\mu) = \mathbf{N}\) 。对每个 \(a \in \mathbf{N}\) 及 \(a\) 的每个邻域 V，存在一个函数 \(f \in \mathcal{K}(\mathbf{X}; \mathbf{C})\) ，其支集含于 V，在点 \(a\) 处等于 1 而在 N 的其他点处等于 0，于是 \(\mu(f) = h(a) \neq 0\) 。另一方面，若 \(b \notin \mathbf{N}\) ，则存在 \(b\) 的一个不与 N 相交的邻域 W；于是对每个支集含于 W 的函数 \(g \in \mathcal{K}(\mathbf{X}; \mathbf{C})\) ，有 \(\mu(g) = 0\) ，这就证明 \(b \notin \operatorname{Supp}(\mu)\) 。

反之，设 \(\mu\) 是使 \(\operatorname{Supp}(\mu)\) 为 X 的离散闭子空间 N 的一个测度。对每个 \(a \in N\) ，存在 a 的一个不含 N 中异于 a 的点的开邻域 \(V_{a}\) ；因此 \(\mu\) 在 \(V_{a}\) 上的限制是支集为 \(\{a\}\) 的点测度（No.~2, 命题 5），从而（No.~4, 命题 12）具有 \(h(a)\varepsilon_{a}\) 的形式，其中 \(h(a)\neq0\) 。令 \(h(x)=0\) 在 \(\complement N\) 的各点处，并以 \(\nu\) 记由诸质量 \(h(x)\) 定义的测度，则局部化原理表明 \(\nu=\mu\) 。

由此可知，在离散空间 X 上，每个测度都是离散的。

